\subsection{Ordinary classes and the three-valuation in the fifty-three factor}

\begin{statement}{E.124}{Proved}{A saturated ordinary-unit basis, the next full class group and the exact bounded valuation image}
Keep the compatible roots and actual ideal maps of E.123. Write
\[
 \begin{gathered}
 L_d=\mathbb Q(\beta_d),\quad
 \beta_d^{3^d}=53\zeta_3,\quad \beta_{d+1}^3=\beta_d,\\
 D_d=\operatorname{Cl}(L_d)\{3\},\quad
 C_d=\operatorname{Cl}_S(L_d)\{3\},\quad
 E_d=\mathcal O_{L_d}^{\times},\qquad S\mid\{3,53\},\\
 q_d=[\mathfrak q_{3,d}],\quad t_d=[\mathfrak q_{53,d}],\quad
 N_d:D_{d+1}\longrightarrow D_d,\quad
 j_d:D_d\longrightarrow D_{d+1}.
 \end{gathered}
 \tag{E124.1}
\]
Here $N_d$ is ideal norm and $j_d$ is ideal extension. The prime
over fifty-three is the principal ideal $(\beta_d)$.
The eight literal ordinary units $\epsilon_1,\ldots,\epsilon_8$
in the first eight columns of E.122's upper witness list, together
with $\zeta_3$, give the following isomorphism after completion
for multiplication by $3^n$:
\[
 \mu_3\times\mathbb Z_3^8\xrightarrow{\sim}\widehat E_{2,3},
 \qquad (z,a_1,\ldots,a_8)\longmapsto z\prod_i\epsilon_i^{a_i}.
 \tag{E124.2}
\]
The inverse is the exponent map in this proved basis; its reduction
modulo three is the explicit inverse residue matrix in the certificate.
Every ordinary unit of $L_2$ is a field norm from $L_3$.

In E.123's actual three-ideal basis and the Smith basis below,
the entire ordinary three-class groups and actual maps are
\[
 \begin{gathered}
 D_2=C_9\oplus C_9\oplus C_3,\qquad
 D_3=C_{27}\oplus C_9\oplus C_3,\\
 q_2=(6,3,1),\qquad q_3=(21,0,1),\\
 N_2=\begin{pmatrix}8&1&0\\1&0&0\\2&1&1\end{pmatrix},
 \qquad
 j_2=\begin{pmatrix}0&3&0\\3&3&0\\0&0&0\end{pmatrix}.
 \end{gathered}
 \tag{E124.3}
\]
Each row is reduced modulo its target order. Thus $q_2$ has order
three and $q_3$ has order nine. The actual automorphism
$\sigma_2(\beta_3)=\zeta_3\beta_3$ acts identically on $D_3$.
The full kernels and quotient maps are
\[
 \begin{gathered}
 \ker N_2=\langle(9,0,0)\rangle=\langle3q_3\rangle,\qquad
 \ker j_2=\langle(3,0,0),(0,0,1)\rangle,\\
 N_2j_2=3I_{D_2},\qquad j_2N_2=3I_{D_3},\\
 P_2=\begin{pmatrix}7&7&0\\2&4&3\end{pmatrix},\qquad
 P_3=P_2N_2=\begin{pmatrix}0&7&0\\8&5&3\end{pmatrix}\pmod9,\\
 P_i:D_i\twoheadrightarrow C_2=C_9c_{31}\oplus C_9c_{103},
 \qquad \ker P_i=\langle q_i\rangle\quad(i=2,3).
 \end{gathered}
 \tag{E124.4}
\]
For $P_3$, the identification $C_3\to C_2$ is E.123's actual
norm isomorphism. The subscripts $31,103$ denote its literal
pivot primes.

The ordinary-unit norm defect and completed kernel are
\[
 \begin{gathered}
 E_2/N_{L_3/L_2}E_3\xleftarrow{\sim}
 D_3/(3D_3+\langle q_3\rangle)\xrightarrow{\sim}C_3^2,\\
 [(r,s,t)]\longmapsto(r\bmod3,s\bmod3),\qquad
 \ker(\widehat E_{3,3}\to\widehat E_{2,3})
                         \simeq\mu_3\oplus\mathbb Z_3^{18}.
 \end{gathered}
 \tag{E124.5}
\]
The first arrow is the actual Hilbert--90 ideal-to-unit map,
with the inverse below. Its defining norm subgroup is retained;
no numerical two-row character on the eight-unit basis is assigned.
The torsion factor in the kernel is the literal root subgroup.

Define actual integers $a_d$ by $\operatorname{ord}(q_d)=3^{a_d}$.
Then
\[
 \begin{gathered}
 a_2=1,\quad a_3=2,\qquad
 a_d\le a_{d+1}\le a_d+1,\qquad h:=\sup_{d\ge2}a_d<\infty,\\
 D_d\simeq C_{3^{a_d+1}}\oplus C_9\oplus C_3\quad(d\ge3),
 \qquad
 \varprojlim_{d,N_d}D_d\simeq C_{3^{h+1}}\oplus C_9\oplus C_3.
 \end{gathered}
 \tag{E124.6}
\]
The proof gives actual compatible basis choices and inverse maps.
In those bases $q_d=(3,0,1)$; for $d\ge3$ each norm reduces the
first coordinate and preserves the other two. Extension multiplies
any norm lift by three; its ambiguity is killed by three.
The finite integer $h\ge2$ is not numerically evaluated here.
Without any basis choice, for every $d\ge2$,
\[
 \begin{gathered}
 N_d\text{ is onto},\quad
 \ker N_d\subset\langle q_{d+1}\rangle,\quad
 |\ker N_d|=3^{a_{d+1}-a_d}\le3,\\
 (\sigma_d-1)^2D_{d+1}=3(\sigma_d-1)D_{d+1}=0,\qquad
 j_dN_d=3I_{D_{d+1}},\quad N_dj_d=3I_{D_d},\\
 E_d/N_dE_{d+1}\simeq C_3^{\,3-a_{d+1}+a_d},\qquad
 \ker(\widehat E_{d+1,3}\to\widehat E_{d,3})
                   \simeq\mu_3\oplus\mathbb Z_3^{\,2\cdot3^d}.
 \end{gathered}
 \tag{E124.7}
\]
The finite unit quotient keeps its actual norm subgroup. The
proof supplies its Hilbert--90 comparison as well as its order.

Let $\mathcal V$ be E.123's coherent completed $S$-unit group and
$\mathcal E=\varprojlim_{d,N}\widehat E_{d,3}$. For its literal
valuation maps one has
\[
 \begin{gathered}
 \nu_3(\mathcal V)=3^h\mathbb Z_3,\qquad
 (\nu_3,\nu_{53})(\mathcal V)=3^h\mathbb Z_3\times2\mathbb Z_3.\\
 \text{Choose }\gamma\in\mathcal V\text{ with }
 (\nu_3,\nu_{53})(\gamma)=(3^h,0),\qquad
 \boldsymbol\beta=(\beta_d)_d.\\
 \mathcal E\times\mathbb Z_3^2\xrightarrow{\sim}\mathcal V,
 \qquad (u,k,t)\longmapsto u\gamma^k\boldsymbol\beta^t,\\
 y\longmapsto
 \left(y\gamma^{-\nu_3(y)/3^h}\boldsymbol\beta^{-\nu_{53}(y)/2},
       \frac{\nu_3(y)}{3^h},\frac{\nu_{53}(y)}2\right).
 \end{gathered}
 \tag{E124.8}
\]
All maps are continuous and both compositions are identities.
Their literal rational norm is
\[
 N_{\mathbb Q}(u\gamma^k\boldsymbol\beta^t)=3^{\,3^hk}53^{\,2t}.
 \tag{E124.9}
\]
Together with E.123's fixed divisor splitting $s$, the entire
positional factor has the explicit coordinates
\[
 (u,k,t,v)\longmapsto u\gamma^k\boldsymbol\beta^t s(v),\qquad
 N_{\mathbb Q}=3^{\,3^hk+\rho(v)}53^{\,2t+\eta(v)}
                         \prod_{p\in\mathcal P}p^{v_p}.
 \tag{E124.10}
\]
The inverse takes the actual outside-$S$ divisor $v$, then applies
(E124.8) to $x/s(v)$. Replacing $\gamma$ by $e\gamma$,
$e\in\mathcal E$, changes $u$ to $ue^{-k}$ and preserves
the valuation coordinates. These are E.46's actual total-norm
coordinates for this positional factor. The remaining numerical
invariant is $h$; the other factors and intrinsic lifting maps
retain their definitions in E.23. The manuscript anchor is
\S51.16, fols.~733v--736v, scans172--178, through E.46's
root-cubing and Kummer--Shapiro maps.
\end{statement}

\begin{proof}
\emph{Literal units and saturation.}
The field $L_2$ has degree eighteen, is totally imaginary and
has ordinary unit rank eight. Its three-primary roots of unity
are exactly $\mu_3$, by E.123's horizontal root calculation.
The rational coefficient vectors of the eight $\epsilon_i$
are retained in the immutable witness input. E.122's independent
relative-ring identities and complete principal divisors show
that they are ordinary units. No regulator or upper class-number
output is used for this assertion.

At a degree-one prime with residues $\beta_2=b$, $\zeta_3=z$,
use the actual character $u\mapsto\log_z(\overline u^{(p-1)/3})$.
The new certificate selects the nine primes
\[
 \begin{array}{c|rrrrrrrrr}
 p&7&7&7&31&31&31&43&79&97\\
 b&1&2&4&11&24&27&4&41&44\\
 z&2&2&2&25&25&25&36&55&35
 \end{array}.
 \tag{E124.11}
\]
They are actual maximal-order primes: these roots of
$X^{18}+53X^9+53^2$ are simple, and the primes avoid its
discriminant support $\{3,53\}$. On the columns
$(\zeta_3,\epsilon_1,\ldots,\epsilon_8)$ the residue matrix is
\[
 F=\begin{pmatrix}
 2&0&0&2&0&1&2&0&0\\
 2&0&0&0&0&0&1&2&2\\
 2&0&0&2&1&1&2&1&1\\
 1&2&1&0&2&0&0&2&1\\
 1&2&1&0&0&1&1&0&2\\
 1&2&1&1&2&2&0&1&0\\
 2&0&0&0&0&0&2&1&1\\
 2&0&2&0&2&0&2&0&1\\
 2&0&1&1&1&1&2&0&1
 \end{pmatrix}\pmod3.
 \tag{E124.12}
\]
The certificate retains $F^{-1}$ and verifies both products.
Since $E_2/E_2^3$ has dimension nine, these nine classes are
its full basis. Their subgroup has rank eight modulo torsion
and index prime to three. Completion gives (E124.2).
One inverse procedure reads the exponents modulo three with
$F^{-1}$, divides by that literal unit product, and takes the
now existing cube in $E_2$; iteration recovers the exponent digits.

At the unique fifty-three place the residue field is
$\mathbb F_{53}[z]/(z^2+z+1)$ with the specified reduction of
$\zeta_3$. The positive Artin character for $L_3/L_2$ on units is
\[
 \chi_{53}(u)=-\log_z(\overline u^{936})\pmod3,\qquad
 936=(53^2-1)/3.
 \tag{E124.13}
\]
This is the tame symbol of $(\beta_2,u)$ in E.119's arithmetic
Artin convention. Its value on every basis element in
(E124.2) is zero, by direct finite-field evaluation. Hence all
ordinary units are local norms there. At an unramified place
every unit is a norm, and the product of the local Artin
characters makes the character at the sole three-place zero.
There are no real places. Milne VIII Theorem3.1, with its
injective Tate-cohomology proof, now gives
$E_2\subset N_{L_3/L_2}L_3^\times$ \cite{E_MilneCFT}.
The local-to-global product is V Proposition5.2 and
Theorem5.3; V Example5.4 states its Hilbert-symbol form.

\emph{The two full ordinary class groups.}
E.119 gives $\operatorname{Cl}(L_1)\{3\}=C_9\oplus C_3$
and $\operatorname{Cl}_S(L_1)\{3\}=C_3^2$.
Because $(\beta_1)=\mathfrak q_{53,1}$, their kernel is
generated by $q_1$ and has order three. E.123's literal
horizontal norm $w$ has complete divisor
$3\mathfrak q_{3,2}-\mathfrak q_{53,2}$.
Thus $w\beta_2$ has divisor $3\mathfrak q_{3,2}$, while
$N_1q_2=q_1$ proves that $q_2$ has exact order three.
Removing the $S$-primes gives the actual exact sequence
\[
 0\longrightarrow\langle q_d\rangle\longrightarrow D_d
       \longrightarrow C_d\longrightarrow0.
 \tag{E124.14}
\]
The class $q_d$ is three-primary: extensions of the relation
for $q_2$ annihilate it by powers of three. The other removed
class is zero, so there is no suppressed prime-to-three term.

At level two, $|C_2|=81$ gives $|D_2|=243$.
E.123's literal ideals $I_1,I_2,I_3$ have complete principal
power relations of orders nine, nine and three. Their classes
generate $C_2$, and their combination $(6,3,1)$ is $q_2$.
They generate $D_2$; the resulting surjection from
$C_9\oplus C_9\oplus C_3$ is an isomorphism by equal orders.
This proves its actual ideal basis in (E124.3).

The cyclic cubic $L_3/L_2$ has exactly two totally ramified
finite places and no infinite ramification. The field-unit
norm index just proved is one. The three-primary ambiguous
class number formula therefore gives
\[
 |D_3^{\sigma_2}|=|D_2|\,3^{2-1}=729.
 \tag{E124.15}
\]
This is Lemmermeyer Theorem1(1), proved by its fixed ideal
calculation and the unit Herbrand quotient
\cite{E_LemmermeyerAmbiguous}. Extension of $3q_2=0$ gives
$9q_3=0$. Equation (E124.14) and $|C_3|=81$ bound
$|D_3|$ by $729$. Equality proves that $q_3$ has order nine
and that $\sigma_2$ fixes the entire ordinary group.

\emph{The actual idelic presentation and local signs.}
Let $H_3$ be the ordinary Hilbert three-class field of $L_3$.
It is Galois over $L_2$. The group
$\Gamma=\operatorname{Gal}(H_3/L_2)$ is an extension of
$D_3$ by a cyclic group with trivial action. Choosing a lift
of that generator proves that $\Gamma$ is abelian, of order $2187$.

The two local unit norm quotients have order three; denote their
positive generators by $e_3,e_{53}$. The literal norm
uniformizers are $1+\beta_2=N(1+\beta_3)$ and
$\beta_2=N\beta_3$. They fix the local valuation splittings
and the zero uniformizer characters. The other three generators
are actual lifts of $I_1,I_2,I_3$. Complete factorization in
the certified maximal order gives their Frobenius characters
$(1,0,2)$ in $L_3/L_2$.

Specifically, $I_1$ is the prime above $2683$ with residues
$(\beta_2,\zeta_3)=(1567,2046)$; $I_2$ is the prime above
$7$ with residues $(1,2)$. The ideal $I_3$ is the product
of that prime, the prime above $31$ with residues $(24,25)$,
and the stated prime above $2683$. Their factor characters
are $1$, $0$, and $(0,1,1)$ respectively. The nf-only
program checks these ideal factorizations and every residue;
in particular it does not treat $I_3$ as prime.

For the retained principal generators $a_i$ of
$I_1^9,I_2^9,I_3^3$, direct use of (E124.13) gives
\[
 \begin{array}{c|ccc}
 &a_1&a_2&a_3\\ \hline
 \chi_3&0&2&0\\
 \chi_{53}&0&1&0
 \end{array}.
 \tag{E124.16}
\]
Their residue pairs in $\mathbb F_{53}[z]/(z^2+z+1)$
are $(8,16),(43,36),(0,45)$. The three-place row follows
from reciprocity: every outside-$S$ exponent in these complete
principal divisors is divisible by three. All global ordinary
unit characters vanish. The exact idelic presentation is consequently
\[
 R=\begin{pmatrix}
 3&0&0&2&0\\0&3&0&1&0\\0&0&9&0&0\\
 0&0&0&9&0\\0&0&0&0&3
 \end{pmatrix},\qquad
 \Gamma=\mathbb Z_3^5/R\mathbb Z_3^5,\qquad
 \phi=(1,1,1,0,2)\pmod3.
 \tag{E124.17}
\]
Here $\phi$ is restriction to the radical $\beta_3$.
The positive power-column sign is fixed by the principal
idele relation
$nI_i+\chi_3(a_i)e_3+\chi_{53}(a_i)e_{53}=0$.
Reducing any idele by its ideal class and then its local
unit components gives a surjection from this presentation
to $\Gamma$. Its determinant is $2187$, so the map is an
isomorphism. These are the actual idele and Artin maps of
Milne V Proposition5.2 and Theorems5.3 and5.5.
Local norm compatibility follows from the proved cup-product
Res/Cor identities in III Theorem3.1: in Tate degree $-2$
the group maps are transfer and inclusion, giving respectively
inclusion/transfer and norm/restriction on Artin maps.
No proof is attributed to III Lemma3.2, whose printed proof
is omitted.

The full preimage of $\ker\phi$ in $\mathbb Z_3^5$ has basis
\[
 K=\begin{pmatrix}
 3&-1&-1&0&-2\\0&1&0&0&0\\0&0&1&0&0\\
 0&0&0&1&0\\0&0&0&0&1
 \end{pmatrix},\qquad \text{relations }K^{-1}R.
 \tag{E124.18}
\]
The certificate retains unimodular matrices $U,V$ satisfying
$U(K^{-1}R)V=\operatorname{diag}(27,9,3,1,1)$.
The actual forward map is $UK^{-1}$ on the kernel lattice;
the inverse is $KU^{-1}$ on Smith representatives. Complete
ambient matrices, including thirds and the trivial Smith
coordinates, are serialized and decoded exactly.

The actual relation
$\mathfrak q_{3,2}=I_1^6I_2^3I_3(x_Q)$ has
$\chi_{53}(x_Q)=2$. Removing its three-valuation with
the norm uniformizer $1+\beta_2$ gives
\[
 \chi_3(x_Q)=6\cdot1+3\cdot0+1\cdot2-2=0\pmod3.
 \tag{E124.19}
\]
Thus the upper three-prime has kernel vector $(0,-2,6,3,1)$.
Applying $UK^{-1}$ gives $(21,0,1,0,0)$.
Projection onto the last three idelic coordinates is the
actual norm to $D_2$; it gives $N_2$ in (E124.3).
The trace identity gives $j_2(b)=3a$ for every lift
$N_2(a)=b$. The norm kernel has exponent three, so this
rule is independent of the lift and gives the displayed $j_2$.
The independent verifier checks the rule on every lift of
each of the $243$ lower classes, both trace compositions,
all $729$ upper classes, both kernels and both $S$-quotient maps.

\emph{Unit defects and completed kernels.}
An invariant ideal is a product of extended lower ideals and
ramified primes: its coefficients are constant on each
unramified orbit, and ramification gives the extra prime
generator. The fifty-three prime is principal. Since $N_2$
is onto and its trace is three, the strongly ambiguous subgroup
is $3D_3+\langle q_3\rangle$. Its order is $81$, and the
two-row quotient in (E124.5) has exactly this kernel.

The Hilbert--90 maps have literal inverses. Given an invariant
class represented by a divisor $A$, choose $x$ with
$(\sigma_2-1)A=\operatorname{div}(x)$ and send it to $Nx$.
This is an ordinary unit, well-defined modulo unit norms.
Conversely, for $u=Nx\in E_2$, the divisor of $x$ has norm
zero and equals $(\sigma_2-1)A$. On a split orbit this is
the difference description of the augmentation kernel; at
a fixed prime the coefficient must be zero. Changing the
norm witness by Hilbert90 changes $A$ by a principal divisor
and an invariant divisor. The latter is exactly the strongly
ambiguous ambiguity. These maps are inverse, as in the
ideal-level proof of Lemmermeyer Proposition1.

A finite-index subgroup of a finitely generated abelian group
has its induced three-power topology cofinal with its own:
an integer killing the quotient gives both inclusions of powers.
Thus the unit norm sequence and its finite quotient remain
exact after completion. The ranks at levels $d,d+1$ are
$3^d-1,3^{d+1}-1$; lower cubes are unit norms, so the image
has finite index and the free kernel rank is $2\cdot3^d$.
The literal root subgroups are $\mu_3$, and norm is cubing,
so the kernel torsion is exactly $\mu_3$.
Choose a basis of the kernel modulo torsion and lifts $k_i$.
The displayed kernel isomorphism is
$(z,(a_i))\mapsto z\prod_i k_i^{a_i}$; its inverse takes
quotient coordinates, divides by this product and reads the
remaining root. This specifies every basis and lift choice.

\emph{Later actual norms and bounded ordinary growth.}
Every $N_d$ is onto. Indeed the ramified cyclic cubic
$L_{d+1}/L_d$ intersects the ordinary Hilbert three-class
field of $L_d$ only in $L_d$. Extend a prescribed Hilbert
Frobenius to the compositum with trivial component on
$L_{d+1}$ and choose a prime with that Frobenius. A relative
degree-one prime above it has the prescribed ideal norm.
This uses the actual norm/Frobenius square of Milne V3.3,
the Hilbert-field identification V3.6/Example3.9, and VIII7.4.

Total ramification gives $N_dq_{d+1}=q_d$ and
$j_dq_d=3q_{d+1}$, proving
$a_d\le a_{d+1}\le a_d+1$.
The $S$-class norm is an isomorphism, so
$\ker N_d\subset\langle q_{d+1}\rangle$ and has order at
most three. Also $N_d\sigma_d=N_d$ and
$\sigma_dq_{d+1}=q_{d+1}$. Consequently $T=\sigma_d-1$
has $T^2=3T=0$. Norm of an ideal followed by extension is
the product of its conjugates, so
$j_dN_d=1+\sigma_d+\sigma_d^2=3+3T+T^2=3$.
The other composition is three as well.

Set $F=L_2$ and
\[
 F_\infty=F(\mu_{3^\infty},53^{1/3^\infty}).
 \tag{E124.20}
\]
Each literal $L_d$ is a subfield, through
$\beta_d=53^{1/3^d}\zeta_{3^{d+1}}$ with the fixed roots.
There is one place of $F$ over three. Kato Theorem1.1(2)
gives finite index for the image of the actual inverse
system of completed three-units over all finite subfields
of (E124.20) in
$\mathcal O_F[1/3]^\times\otimes\mathbb Z_3$
\cite{E_KatoUniversalNorms}. Restriction to the coordinates
$L_d$ is an actual map of norm systems. Relative residue
degree one makes the three-valuation unchanged. The
finite-index image at $F$ has nonzero three-valuation, since
the whole valuation image there is $3\mathbb Z_3$.
That nonzero value is divisible by every $3^{a_d}$.
Hence $h<\infty$. Kato supplies boundedness and coherent
systems here, not a numerical value of its finite index.

\emph{Higher bases and inverse maps.}
The nondecreasing bounded integers $a_d$ are eventually
constant. Surjectivity and equal orders then make the
ordinary norms isomorphisms. Their inverse limit is that
eventual finite group and maps onto $D_3$. Its coherent
prime class $q$ has order $3^h$. In $D_3$ put
$x_3=(7,0,0)$ and $y_3=(0,1,0)$; then
$9x_3=3q_3$ and $9y_3=0$.
Choose actual norm lifts $x,y$ in the inverse group.
Their $S$-quotient coordinates are a basis. Thus
\[
 9x=3u q,\qquad9y=3v q,\qquad
 u\equiv1\pmod3,\quad v\equiv0\pmod3.
 \tag{E124.21}
\]
Indeed their coefficients on $q$ reduce modulo nine to
three and zero under the actual projection to $D_3$.
Choose representatives for $u,v$, and define
\[
 X=u^{-1}x,\qquad Y=y-vX,\qquad Z=q-3X.
 \tag{E124.22}
\]
Then $9X=3q$, $9Y=0$, $3Z=0$ and $q=3X+Z$.
The three resulting cyclic groups surject onto the inverse
class group and have total order $81\cdot3^h$, its order
by (E124.14). Hence this is an isomorphism.
An inverse reads quotient coordinates $r,s$ modulo nine
in $\overline X,\overline Y$, writes the residual class
as $tq$, and returns $(r+3t,s,t\bmod3)$.
Changing $r$ by nine changes $t$ by minus three, so the
answer is unchanged. The other ambiguities are killed
by $9Y=0$ and the order of $q$.

Taking actual norms of $X,Y,Z$ gives bases at every
$d\ge3$ with the same formulas and $h$ replaced by $a_d$.
The inverse construction commutes with norm, which proves
the coordinate maps in (E124.6). Its extension rule follows
from the trace identity and the norm kernel's exponent three.
For all $d\ge2$ the strongly ambiguous subgroup is therefore
$3D_{d+1}+\langle q_{d+1}\rangle$ and has index nine.
Lemmermeyer Theorem1(2) gives
\[
 [E_d:N_dE_{d+1}]
   =\frac{3|D_d|}{|D_{d+1}|/9}
   =3^{\,3-a_{d+1}+a_d}.
 \tag{E124.23}
\]
All lower cubes are unit norms, proving the asserted
elementary group type. The actual subgroup is still
$N_dE_{d+1}$. For its field-norm part the same Hilbert--90
map uses $D_{d+1}^{\sigma_d}$ and quotients by the strongly
ambiguous subgroup. This specifies the comparison even
where the individual $\sigma_d$ action has not been evaluated.
The earlier rank and literal-root computation proves the
remaining completed-kernel claim.

\emph{The coherent valuation image.}
Put $W_d=\mathcal O_{L_d}[1/3]^\times$. Its completed
three-valuation image is exactly $3^{a_d}\mathbb Z_3$,
by (E124.14). For
$e\in\bigcap_d3^{a_d}\mathbb Z_3$, every completed valuation
fibre at level $d$ is nonempty and compact. A finite chain
of norm equations is solved by choosing in its highest
fibre and taking norms. In the product of compact fibres,
these closed equations have the finite intersection
property. Their intersection proves both directions of
\[
 v_3\!\left(\varprojlim_{d,N}\widehat W_{d,3}\right)
     =\bigcap_d3^{a_d}\mathbb Z_3=3^h\mathbb Z_3.
 \tag{E124.24}
\]
Choose $\gamma$ in this limit with valuation $3^h$.
The literal tuple $\boldsymbol\beta$ is norm coherent,
has zero three-valuation and fifty-three valuation one.
In any coherent $S$-unit tuple the latter valuation is
constant. Divide by the indicated power of
$\boldsymbol\beta$, then by the indicated power of
$\gamma$; the result is a coherent ordinary unit.
This proves (E124.8) and both compositions. Compact
finite-level unit groups and continuous three-adic
exponentiation also prove continuity of each inverse.

An ordinary unit has absolute norm $\pm1$, which is
trivial in the three-power completion of $\mathbb Q^\times$.
The two absolute residue degrees are one and two.
These give (E124.9) on the actual norms and then
(E124.10) using E.123's fixed divisor splitting.
The numerical value of $h$ and a literal coherent
minimal-valuation tuple remain the specified next
arithmetic step. All stage-three computations, boundedness,
the exact image and the displayed comparison maps are
proved without assigning that value.
\end{proof}
